\documentclass[12pt, a4paper]{article}

% ── Packages ──────────────────────────────────────────────────────────────────
\usepackage[margin=2.5cm]{geometry}
\usepackage{graphicx}
\usepackage{amsmath}
\usepackage{amssymb}
\usepackage{hyperref}
\usepackage{xcolor}
\usepackage{booktabs}
\usepackage{longtable}
\usepackage{array}
\usepackage{enumitem}
\usepackage{fancyhdr}
\usepackage{titlesec}
\usepackage{parskip}
\usepackage{caption}
\usepackage{subcaption}
\usepackage{listings}
\usepackage{mdframed}
\usepackage{tcolorbox}
\usepackage{float}

% ── Colours ───────────────────────────────────────────────────────────────────
\definecolor{darkblue}{RGB}{26, 57, 107}
\definecolor{medblue}{RGB}{47, 84, 150}
\definecolor{lightblue}{RGB}{189, 215, 238}
\definecolor{darkgreen}{RGB}{55, 86, 35}
\definecolor{lightgray}{RGB}{242, 242, 242}
\definecolor{boxborder}{RGB}{47, 84, 150}
\definecolor{warnorange}{RGB}{197, 90, 17}

% ── Section formatting ────────────────────────────────────────────────────────
\titleformat{\section}{\large\bfseries\color{darkblue}}{}{0em}{}[\titlerule]
\titleformat{\subsection}{\normalsize\bfseries\color{medblue}}{}{0em}{}
\titleformat{\subsubsection}{\normalsize\bfseries\color{darkgreen}}{}{0em}{}

% ── Header / Footer ───────────────────────────────────────────────────────────
\pagestyle{fancy}
\fancyhf{}
\fancyhead[L]{\small\color{medblue}Sand Mining Detection System — Technical Documentation}
\fancyhead[R]{\small\color{medblue}Damodar Basin Pilot Study}
\fancyfoot[C]{\thepage}
\renewcommand{\headrulewidth}{0.4pt}

% ── Box environments ──────────────────────────────────────────────────────────
\tcbuselibrary{skins, breakable}

\newtcolorbox{keybox}[1]{
  colback=lightblue!30,
  colframe=boxborder,
  title={\textbf{#1}},
  fonttitle=\small\bfseries,
  breakable
}

\newtcolorbox{notebox}{
  colback=yellow!10,
  colframe=warnorange,
  breakable
}

% ── Listings style ────────────────────────────────────────────────────────────
\lstset{
  basicstyle=\small\ttfamily,
  backgroundcolor=\color{lightgray},
  frame=single,
  rulecolor=\color{medblue},
  breaklines=true,
  captionpos=b
}

% ─────────────────────────────────────────────────────────────────────────────
\begin{document}

% ── Title Page ────────────────────────────────────────────────────────────────
\begin{titlepage}
  \centering
  \vspace*{2cm}

  {\Huge\bfseries\color{darkblue}
    Automated Sand Mining Detection\\[0.4em]
    Using Satellite Imagery and\\[0.4em]
    Machine Learning
  \par}

  \vspace{0.5cm}
  \rule{12cm}{1.5pt}
  \vspace{0.5cm}

  {\large\color{medblue}
    A Plain-Language Technical Documentation\\
    of the Sand Mining Detection Pipeline
  \par}

  \vspace{1cm}

  {\normalsize
    \textbf{Pilot Study Area:} Damodar River Basin, West Bengal / Jharkhand, India\\[0.3em]
    \textbf{Primary Sensor:} Sentinel-2 Multispectral (10\,m resolution)\\[0.3em]
    \textbf{Temporal Coverage:} 2023--2026 (3--5 year historical window)\\[0.3em]
    \textbf{Mapping Scale:} 0.4--0.5\,km point spacing along river corridor\\[0.3em]
    \textbf{Ultimate Goal:} All-India multi-basin sand mining probability map
  \par}

  \vspace{2cm}

  \begin{tcolorbox}[colback=lightblue!20, colframe=darkblue, width=12cm]
    \centering
    \textbf{What This System Does in One Sentence}\\[0.5em]
    \normalsize
    It automatically looks at satellite photos of Indian rivers taken over
    the past 3--5 years, detects visual and statistical patterns associated
    with illegal sand mining, and produces a map showing exactly where
    mining is likely occurring right now --- without anyone having to
    physically visit those locations.
  \end{tcolorbox}

  \vfill
  {\small March 2026}
\end{titlepage}

% ── Table of Contents ─────────────────────────────────────────────────────────
\tableofcontents
\newpage

% =============================================================================
\section{Introduction: The Problem This Solves}
% =============================================================================

\subsection{What is Riverine Sand Mining?}

Sand and gravel extracted from riverbeds are the world's most consumed
solid materials by weight. In India alone, construction demand exceeds
500 million tonnes of river sand per year. The vast majority of this
extraction is \textbf{illegal} --- occurring without permits,
environmental clearance, or any monitoring whatsoever.

Sand mining destroys riverbeds. It lowers the riverbed level (a process
called ``channel incision''), causes banks to collapse, cuts off
groundwater recharge, kills fish and river dolphins, and can even cause
bridges to collapse. Despite a legal framework (Sustainable Sand Mining
Guidelines 2016, revised 2020), enforcement is almost non-existent
because no one knows \emph{where} mining is happening.

\subsection{Why Can't We Just Drive Around and Check?}

India has hundreds of thousands of kilometres of rivers. The Damodar
River alone --- the pilot study area in this system --- runs
361\,km through West Bengal and Jharkhand. Physical inspection of every
stretch is:

\begin{itemize}[noitemsep]
  \item \textbf{Impossible at scale} --- too many rivers, too few inspectors
  \item \textbf{Dangerous} --- sand mining mafias are known to be violent
  \item \textbf{Too slow} --- illegal operations move quickly
  \item \textbf{Expensive} --- field surveys cost lakh of rupees per campaign
\end{itemize}

\subsection{What This System Does Instead}

Satellites photograph every point on Earth's surface every 5--10 days
for free. This system:

\begin{enumerate}[noitemsep]
  \item Downloads satellite images of the river
  \item Extracts mathematical features from each image (colour patterns,
        texture, vegetation cover, water turbidity, and how these have
        changed over 3--5 years)
  \item Uses a machine learning model trained by a human expert to classify
        each location as ``Sand Mining Likely'', ``Possible'', or ``No Mining''
  \item Produces an interactive HTML map showing the results
\end{enumerate}

% =============================================================================
\section{The Satellite Data: What We See from Space}
% =============================================================================

\subsection{Primary Sensor: Sentinel-2}

\begin{keybox}{Sentinel-2 at a Glance}
  \begin{itemize}[noitemsep]
    \item \textbf{Operated by:} European Space Agency (ESA)
    \item \textbf{Spatial resolution:} 10 metres per pixel (visible/NIR bands)
    \item \textbf{Revisit time:} Every 5 days at the equator
    \item \textbf{Data access:} Free and open via Google Earth Engine
    \item \textbf{Available since:} 2017 (harmonised archive)
    \item \textbf{Bands used:} 10 spectral bands from Blue to SWIR
  \end{itemize}
\end{keybox}

At 10\,m resolution, each pixel represents a 10\,m $\times$ 10\,m patch
on the ground. This is fine enough to see:
\begin{itemize}[noitemsep]
  \item Individual sand barges on a river
  \item Exposed sandy excavation pits
  \item Disturbed bank vegetation
  \item Turbid (muddy) water plumes from active dredging
\end{itemize}

\subsection{Backup Sensors}

When Sentinel-2 imagery is unavailable (e.g., heavy cloud cover), the
system automatically falls back to:

\begin{table}[H]
\centering
\begin{tabular}{llll}
\toprule
\textbf{Sensor} & \textbf{Resolution} & \textbf{Archive Since} & \textbf{Use} \\
\midrule
Sentinel-2 & 10\,m & 2017 & Primary \\
Landsat 8/9 & 30\,m & 2013/2021 & Fallback (current) \\
Landsat 7 & 30\,m & 1999 & Historical fallback \\
Landsat 5 & 30\,m & 1984 & Deep historical \\
VIIRS & 375\,m & 2012 & Turbidity / SSC proxy \\
\bottomrule
\end{tabular}
\caption{Satellite sensors used in the detection pipeline}
\end{table}

\subsection{Cloud Cover Handling}

India's monsoon season (June--September) covers most river areas with
cloud 60--90\% of the time. The system handles this by:
\begin{itemize}[noitemsep]
  \item Filtering images with $<35\%$ cloud cover (preferred)
  \item Relaxing to $<60\%$ if nothing better is available
  \item Using quarterly composites across 3-month windows to maximise
        the chance of finding a clear image in each period
\end{itemize}

\begin{notebox}
  \textbf{Known limitation:} Sentinel-1 SAR (radar) imagery that can
  see through clouds is not yet integrated. This is the next planned
  upgrade for monsoon-season detection.
\end{notebox}

% =============================================================================
\section{Spectral Indices: Turning Pixels into Information}
% =============================================================================

A satellite image records how much light is reflected from the ground
in different wavelengths (colours). By combining these bands
mathematically, we create ``indices'' --- single numbers that capture
specific ground conditions. Sand mining disturbs the landscape in ways
that show up strongly in these indices.

\subsection{The Six Core Indices}

\subsubsection{NDVI --- Normalised Difference Vegetation Index}

\[
  \text{NDVI} = \frac{\text{NIR} - \text{Red}}{\text{NIR} + \text{Red}}
\]

Measures how much green vegetation is present. Values range from $-1$
(no vegetation / bare soil / water) to $+1$ (dense vegetation).

\textbf{Why it matters for mining:} Active mining sites have vegetation
stripped away. Over time, a legitimate forest or riverside field will
show steadily positive NDVI. A newly mined area will show a sudden
sharp drop in NDVI.

\subsubsection{NDWI --- Normalised Difference Water Index}

\[
  \text{NDWI} = \frac{\text{Green} - \text{NIR}}{\text{Green} + \text{NIR}}
\]

Detects open water bodies. High values ($>0$) indicate water.

\textbf{Why it matters:} Helps separate in-river dredging from
bankside excavation. Also helps distinguish natural sandbars from
active extraction pits.

\subsubsection{MNDWI --- Modified NDWI}

\[
  \text{MNDWI} = \frac{\text{Green} - \text{SWIR1}}{\text{Green} + \text{SWIR1}}
\]

Better at detecting turbid (sediment-laden) water than plain NDWI.
This is particularly useful in Indian rivers which carry high
sediment loads.

\textbf{Why it matters:} Active dredging creates muddy water plumes
downstream. MNDWI detects these turbidity anomalies.

\subsubsection{BSI --- Bare Soil Index}

\[
  \text{BSI} = \frac{(\text{SWIR1} + \text{Red}) - (\text{NIR} + \text{Blue})}
                     {(\text{SWIR1} + \text{Red}) + (\text{NIR} + \text{Blue})}
\]

Detects freshly exposed, bare, dry soil or sediment. Mining sites have
high BSI because:
\begin{itemize}[noitemsep]
  \item Topsoil and vegetation are removed
  \item Sand and gravel have high SWIR and Red reflectance
  \item No moisture-retaining vegetation cover remains
\end{itemize}

\subsubsection{NDBI --- Normalised Difference Built-up Index}

\[
  \text{NDBI} = \frac{\text{SWIR1} - \text{NIR}}{\text{SWIR1} + \text{NIR}}
\]

Originally designed to detect urban areas but also sensitive to
compacted surfaces like equipment staging areas, access roads, and
dried sand stockpiles.

\subsubsection{NDTI --- Normalised Difference Turbidity Index}

\[
  \text{NDTI} = \frac{\text{SWIR1} - \text{SWIR2}}{\text{SWIR1} + \text{SWIR2}}
\]

Proxy for water turbidity (sediment load). Elevated NDTI values in
a river stretch indicate disturbed sediment --- consistent with
active dredging operations.

% =============================================================================
\section{The Historical Feature Innovation}
% =============================================================================

This is the most important technical contribution of this system.

\subsection{The Core Idea}

A single satellite photo cannot reliably tell us if sand mining is
happening. A sandy riverbank and an active mining site can look similar
in one image. But they evolve \emph{differently over time}:

\begin{table}[H]
\centering
\begin{tabular}{lll}
\toprule
\textbf{Feature} & \textbf{Natural Sandbar} & \textbf{Active Mining Site} \\
\midrule
NDVI over 3 years & Stable or seasonal & Declining (vegetation removed) \\
BSI over 3 years & Stable & Increasing (more bare soil) \\
MNDWI over 3 years & Seasonal variation & Elevated turbidity year-round \\
NDWI over 3 years & Follows monsoon & Anomalous during dry season \\
\bottomrule
\end{tabular}
\caption{How natural features vs. mining sites differ over time}
\end{table}

\subsection{How the Temporal Features Are Computed}

For every analysis point, the system:

\begin{enumerate}
  \item Fetches \textbf{12 quarterly Sentinel-2 images} spanning 3 years
        (one image per 3-month window)
  \item For each image, computes the mean NDVI, BSI, MNDWI, and NDWI
        within a 1,500\,m buffer around the point
  \item Fits a \textbf{linear regression line} through the 12 values over time
  \item Extracts the \textbf{slope} of that line as a feature
\end{enumerate}

\begin{keybox}{What the Slopes Mean}
  \begin{itemize}[noitemsep]
    \item \textbf{NDVI\_trend $< 0$:} Vegetation declining over time
          $\rightarrow$ suspicious
    \item \textbf{BSI\_trend $> 0$:} Bare soil increasing over time
          $\rightarrow$ suspicious
    \item \textbf{MNDWI\_trend anomalous:} Water turbidity changing
          unexpectedly $\rightarrow$ suspicious
    \item \textbf{All trends stable:} Natural undisturbed location
  \end{itemize}
\end{keybox}

This approach is validated by:
\begin{itemize}[noitemsep]
  \item Li et al. (2024) --- quarterly BSI/NDWI slopes on the Mekong
        River achieved F1 = 0.85
  \item Bendixen et al. (2021) --- temporal Landsat trends revealed
        20--50\% more sand extraction events than single-date analysis
\end{itemize}

% =============================================================================
\section{Image-Level Features: What the Photo Looks Like}
% =============================================================================

Beyond spectral indices, the system extracts computer vision features
from the visual appearance of each satellite image chip.

\subsection{Colour Statistics}

For each of the Red, Green, Blue, and Grayscale channels:
\begin{itemize}[noitemsep]
  \item Mean brightness
  \item Standard deviation (contrast within the image)
  \item Median value
\end{itemize}

Sand mining sites tend to be brighter (high reflectance from exposed
sand) and more uniform in colour than vegetated natural areas.

\subsection{GLCM Texture Features}

GLCM stands for \textbf{Grey Level Co-occurrence Matrix}. It analyses
how pixel intensities are spatially related to their neighbours ---
essentially measuring the ``texture'' of the image.

Features extracted:
\begin{itemize}[noitemsep]
  \item \textbf{Contrast:} How different adjacent pixels are
  \item \textbf{Homogeneity:} How similar adjacent pixels are
  \item \textbf{Energy:} Uniformity of the texture
  \item \textbf{Correlation:} How linearly related adjacent pixels are
  \item \textbf{Dissimilarity:} Average difference between pixel pairs
  \item \textbf{ASM (Angular Second Moment):} Orderliness of the texture
\end{itemize}

Sand surfaces have low contrast and high homogeneity. Active mining
sites often show high contrast between sand, water, and disturbed soil.

\subsection{Local Binary Pattern (LBP)}

LBP encodes the local texture around each pixel by comparing it to
its eight neighbours. The resulting histogram captures fine-grained
surface texture differences invisible to the naked eye.

\subsection{Edge Detection}

Sobel edge detection identifies sharp boundaries in the image. Mining
sites often have:
\begin{itemize}[noitemsep]
  \item Sharp straight edges (cut banks, access roads)
  \item High edge density near equipment and extraction pits
\end{itemize}

% =============================================================================
\section{The Machine Learning Model}
% =============================================================================

\subsection{What the Model Does}

The model takes all the features described above --- 46 features per
point --- and outputs a single probability between 0 and 1:

\begin{itemize}[noitemsep]
  \item \textbf{0.0--0.5:} No sand mining likely (shown in green on the map)
  \item \textbf{0.5--0.7:} Possible sand mining (shown in orange)
  \item \textbf{0.7--1.0:} Sand mining likely (shown in red)
\end{itemize}

\subsection{Training the Model}

The model learns from \textbf{examples provided by a human expert}.
A domain expert looks at satellite image chips and labels each one:
\begin{itemize}[noitemsep]
  \item \textbf{Label 1:} ``Yes, this looks like sand mining''
  \item \textbf{Label 0:} ``No, this is a natural undisturbed area''
\end{itemize}

The system includes a dedicated \textbf{graphical labeling interface}
(the GUI) that displays each image and lets the expert click a button
to assign the label. Labels are saved to a CSV file for reuse.

\subsection{Six Models Are Trained and Compared}

To ensure the best possible model is selected, the system trains and
evaluates all of the following algorithms on the same data:

\begin{table}[H]
\centering
\begin{tabular}{lp{8cm}}
\toprule
\textbf{Model} & \textbf{How It Works (Plain Language)} \\
\midrule
Random Forest & Many decision trees vote together; majority wins \\
Gradient Boosting & Trees built sequentially, each correcting the previous \\
XGBoost & Optimised gradient boosting, very fast \\
LightGBM & Gradient boosting optimised for large datasets \\
SVM & Finds the widest possible boundary between classes \\
Logistic Regression & Simple linear model of probability \\
\bottomrule
\end{tabular}
\caption{Machine learning models compared in the pipeline}
\end{table}

\subsection{How the Best Model Is Chosen}

The system uses \textbf{5-fold cross-validation} --- it splits the
training data into 5 parts, trains on 4 parts and tests on 1 part,
rotating through all 5 combinations. The model with the best average
F1 score across all 5 folds is selected.

\begin{keybox}{Why F1 Score, Not Just Accuracy?}
  If only 10\% of locations have sand mining, a model that always says
  ``no mining'' would be 90\% accurate but completely useless. F1 score
  penalises both missing real mining sites (false negatives) and false
  alarms (false positives), making it the right metric for this
  imbalanced problem.
\end{keybox}

The system \textbf{prefers ensemble models} (Random Forest, XGBoost,
LightGBM) over SVM when scores are equal, because ensemble models
generalise better to new locations they have not seen during training
--- particularly important with small training datasets.

\subsection{Current Model Performance (Damodar Pilot)}

\begin{table}[H]
\centering
\begin{tabular}{lccccc}
\toprule
\textbf{Model} & \textbf{Accuracy} & \textbf{Precision} & \textbf{Recall} & \textbf{F1} & \textbf{CV F1} \\
\midrule
Random Forest & 0.857 & 1.000 & 0.75 & 0.857 & \textbf{0.820} \\
XGBoost & 0.857 & 1.000 & 0.75 & 0.857 & 0.820 \\
Gradient Boosting & 0.714 & 1.000 & 0.50 & 0.667 & 0.751 \\
LightGBM & 0.571 & 0.571 & 1.00 & 0.727 & 0.714 \\
SVM & 0.714 & 1.000 & 0.50 & 0.667 & 0.820 \\
Logistic Regression & 0.857 & 1.000 & 0.75 & 0.857 & 0.880 \\
\bottomrule
\end{tabular}
\caption{Model comparison results on 25 labeled training samples (Damodar pilot)}
\end{table}

\begin{notebox}
  \textbf{Important caveat:} These results are based on only 25 labeled
  samples. Reliable performance estimates require 100--200+ samples.
  The current results should be treated as preliminary. The high CV F1
  standard deviation ($\pm 0.18$ for RF) reflects this uncertainty.
\end{notebox}

% =============================================================================
\section{The Complete Pipeline: Step by Step}
% =============================================================================

\subsection{Overview}

\begin{figure}[H]
\centering
\begin{tcolorbox}[colback=lightblue!15, colframe=darkblue, width=\textwidth]
\centering
\large
\textbf{STEP 1} \quad $\longrightarrow$ \quad
\textbf{STEP 2} \quad $\longrightarrow$ \quad
\textbf{STEP 3} \quad $\longrightarrow$ \quad
\textbf{STEP 4} \quad $\longrightarrow$ \quad
\textbf{STEP 5}

\normalsize\vspace{0.3em}
Download Images \quad Label Images \quad Train Model \quad Map River \quad View Results
\end{tcolorbox}
\caption{The five-stage pipeline from raw shapefile to interactive map}
\end{figure}

\subsection{Step 1: Generating Sample Points from a Shapefile}

The user provides a \textbf{river shapefile} --- a GIS file that
defines the path of the river as a line. The system:

\begin{enumerate}[noitemsep]
  \item Loads the shapefile using GeoPandas
  \item Reprojects to WGS84 (standard lat/lon coordinates)
  \item Walks along the river line at fixed intervals (e.g., 0.4\,km)
  \item Generates a list of (latitude, longitude) coordinate pairs
\end{enumerate}

For the Damodar River (361\,km) at 0.4\,km spacing, this produces
approximately \textbf{904 sample points}. For training, 30 of these
are randomly selected for image download.

\subsection{Step 2: Downloading Satellite Images}

For each training point, the system contacts \textbf{Google Earth Engine}
and:

\begin{enumerate}[noitemsep]
  \item Creates a 1,500\,m radius circle around the point
  \item Searches for the least-cloudy Sentinel-2 image within the
        past 12 months over that area
  \item Downloads a $512 \times 512$ pixel RGB image as a PNG file
  \item Applies image enhancement (contrast, sharpness, saturation)
        to make visual features clearer for labeling
\end{enumerate}

Each image covers approximately $3 \times 3$\,km on the ground and
takes 4--6 seconds to download.

\subsection{Step 3: Expert Labeling via the GUI}

The graphical labeling interface opens automatically and displays
each downloaded image one by one. The expert:

\begin{itemize}[noitemsep]
  \item Presses \texttt{1} for ``Sand Mining'' or \texttt{0} for ``No Mining''
  \item The image automatically advances to the next
  \item Can press \texttt{?} to skip uncertain images
  \item Presses \texttt{ESC} when done --- labels are saved automatically
\end{itemize}

\textbf{What to look for when labeling:}
\begin{itemize}[noitemsep]
  \item Visible excavation pits or craters near the river
  \item Heavy machinery (excavators, bulldozers) on the riverbank
  \item Barges or boats loaded with sand in the channel
  \item Disturbed, irregular bare soil patches adjacent to water
  \item Unusually turbid or discoloured water
  \item Access roads cut through vegetation to the riverbank
\end{itemize}

\subsection{Step 4: Feature Extraction with Historical Data}

For each labeled image, the system extracts \textbf{46 features}:

\begin{table}[H]
\centering
\begin{tabular}{lll}
\toprule
\textbf{Feature Group} & \textbf{Count} & \textbf{Source} \\
\midrule
Global colour statistics (mean, std, median) & 12 & Current image \\
GLCM texture features & 12 & Current image \\
LBP texture features & 3 & Current image \\
Edge detection features & 3 & Current image \\
Colour ratio features & 4 & Current image \\
Area annotation features & 6 & Expert highlights \\
Historical trend slopes & 4 & 3--5 year GEE time series \\
Historical period count & 1 & GEE \\
Other derived features & 1 & Computed \\
\midrule
\textbf{Total} & \textbf{46} & \\
\bottomrule
\end{tabular}
\caption{Feature breakdown used for model training and prediction}
\end{table}

The \textbf{historical trend extraction} is the most time-consuming
step --- approximately 40 seconds per training point --- because it
requires fetching and processing 12 separate satellite images from
the Earth Engine archive.

\subsection{Step 5: Probability Mapping}

After training, the model is applied to the full river at fine
resolution (0.4--0.5\,km spacing = 723--904 points). For each point:

\begin{enumerate}[noitemsep]
  \item A new current-date satellite image is downloaded
  \item All 46 features are extracted (including 3-year history)
  \item The trained model predicts the probability of sand mining
  \item Results are saved to a CSV file
\end{enumerate}

At 0.4\,km spacing, mapping 904 points with historical analysis takes
approximately \textbf{10--12 hours} (about 45 seconds per point due
to the historical GEE queries).

\subsection{Step 6: Output Generation}

Two outputs are produced automatically:

\begin{enumerate}
  \item \textbf{Interactive HTML Map} --- an OpenStreetMap-based web map
        with colour-coded circles at each analysis point. Clicking any
        circle shows the exact probability, image date, and historical
        change value. The map can be opened in any web browser.

  \item \textbf{Summary PNG Plot} --- a four-panel figure showing:
        the probability distribution histogram, the classification
        counts bar chart, the spatial hotspot heatmap, and the top
        contributing features.
\end{enumerate}

% =============================================================================
\section{Resolution and Coverage}
% =============================================================================

\subsection{Spatial Resolution Hierarchy}

\begin{table}[H]
\centering
\begin{tabular}{lllll}
\toprule
\textbf{Mode} & \textbf{Spacing} & \textbf{Points (Damodar)} & \textbf{Time} & \textbf{Use} \\
\midrule
Quick test & 50\,km & 9 & 8 min & Debugging \\
Coarse & 10\,km & 36 & 30 min & Initial exploration \\
Medium & 5\,km & 72 & 1 hr & Quality check \\
Standard & 1\,km & 361 & 4.5 hr & Regular monitoring \\
Fine & 0.5\,km & 723 & 9 hr & Publication-quality \\
High-res & 0.4\,km & 904 & 11 hr & Maximum detail \\
\bottomrule
\end{tabular}
\caption{Analysis resolution options and approximate runtimes for the Damodar River (361\,km)}
\end{table}

\subsection{Image Chip Coverage}

Each analysis point uses a \textbf{1,500\,m buffer} (3\,km diameter).
At 0.5\,km point spacing, adjacent chips overlap by 1\,km. This means
every location along the river is seen by at least 3 overlapping
analysis windows.

\subsection{Temporal Resolution}

The system looks back \textbf{3--5 years} in history, sampling one
image every \textbf{3 months} (quarterly). For a 3-year window this
gives 12 historical data points per location. Key temporal parameters:

\begin{itemize}[noitemsep]
  \item \textbf{Current image:} Most recent cloud-free image (typically
        within 2 weeks of run date)
  \item \textbf{Historical window:} 3--5 years back (configurable)
  \item \textbf{Historical interval:} 3 months (quarterly)
  \item \textbf{Minimum usable periods:} 3 (system skips trend calculation
        if fewer are available due to cloud cover)
\end{itemize}

% =============================================================================
\section{Scaling to All-India}
% =============================================================================

\subsection{Current Status}

The pipeline has been designed from the start to scale to national
coverage. The Damodar Basin pilot demonstrates the full technical
stack. Scaling to All-India requires:

\begin{enumerate}
  \item \textbf{More labeled training data:} Currently 25 samples.
        A national model needs 500--1,000 samples spanning different
        river types (Himalayan, Deccan, coastal, monsoon-fed)
  \item \textbf{State-wise processing loop:} Process one state at a
        time to avoid Google Earth Engine memory limits
  \item \textbf{River buffer masking:} Use HydroSHEDS or CWC river
        network to restrict analysis to a 2--3\,km riparian corridor
        (reduces compute by $\sim$95\%)
  \item \textbf{Post-processing:} Convert probability points to
        hotspot polygons attributed to district and river name
\end{enumerate}

\subsection{Estimated National Scale}

\begin{table}[H]
\centering
\begin{tabular}{ll}
\toprule
\textbf{Parameter} & \textbf{Estimate} \\
\midrule
Major Indian rivers to cover & $\sim$400 \\
Total river length & $\sim$200,000\,km \\
Points at 0.5\,km spacing & $\sim$400,000 points \\
Processing time (parallel) & 3--5 days \\
GEE quota required & High Volume API \\
Output map resolution & 0.5\,km point spacing, 10\,m image resolution \\
\bottomrule
\end{tabular}
\caption{Estimated parameters for All-India national mapping}
\end{table}

% =============================================================================
\section{Known Limitations and Future Work}
% =============================================================================

\subsection{Current Limitations}

\begin{enumerate}
  \item \textbf{Small training dataset (25 samples):} Reliable
        performance requires 100--200+ labels. The current model
        should be considered a proof-of-concept.

  \item \textbf{No monsoon-season detection:} Sentinel-2 is blocked
        by clouds June--September. Sentinel-1 SAR integration is
        planned but not yet implemented.

  \item \textbf{Only linear trends:} The historical feature uses a
        simple slope (linear regression). Non-linear patterns ---
        such as episodic mining bursts or sudden onset events ---
        are not captured. BFAST or LSTM time-series models would
        improve this.

  \item \textbf{No independent validation:} The model has not yet
        been validated against independently verified ground-truth
        data (e.g., confirmed mining locations from news reports,
        court cases, or field visits).

  \item \textbf{SMOTE not yet installed:} Class imbalance correction
        via SMOTE oversampling is coded but requires \texttt{pip install
        imbalanced-learn}.
\end{enumerate}

\subsection{Planned Upgrades}

\begin{itemize}[noitemsep]
  \item Sentinel-1 SAR fusion for cloud-robust detection
  \item Active learning loop to maximise labeling efficiency
  \item BFAST/LandTrendr breakpoint detection for onset timing
  \item Foundation model features (Prithvi, SatMAE) for small
        dataset performance
  \item State-wise processing loop for All-India scaling
  \item Validation against NGT case locations and news archives
\end{itemize}

% =============================================================================
\section{How to Run the System}
% =============================================================================

\subsection{Prerequisites}

\begin{lstlisting}[language=bash, caption=Installation]
# Python 3.10+, Google Earth Engine account required
pip install earthengine-api geopandas pandas numpy
pip install scikit-learn xgboost lightgbm scipy
pip install Pillow requests tqdm matplotlib seaborn
pip install imbalanced-learn  # for SMOTE
\end{lstlisting}

\subsection{Authentication}

\begin{lstlisting}[language=bash, caption=Google Earth Engine authentication]
gcloud auth application-default login
gcloud config set project YOUR_PROJECT_ID
\end{lstlisting}

\subsection{Running the Full Pipeline}

\begin{lstlisting}[language=bash, caption=Full pipeline command]
python scripts/run_pipeline.py \
  --mode both \
  --shapefile data/raw/your_river.shp \
  --distance 0.5 \
  --years-back 3 \
  --multiple-models \
  --sample-size 100
\end{lstlisting}

\subsection{Key Arguments}

\begin{table}[H]
\centering
\begin{tabular}{lll}
\toprule
\textbf{Argument} & \textbf{Default} & \textbf{Effect} \\
\midrule
\texttt{--mode} & required & \texttt{both / train / map / label} \\
\texttt{--shapefile} & required & Path to river .shp file \\
\texttt{--distance} & 0.2\,km & Point spacing in km \\
\texttt{--years-back} & 0 & Historical window in years \\
\texttt{--sample-size} & 30 & Training images to download \\
\texttt{--multiple-models} & off & Compare all 6 algorithms \\
\bottomrule
\end{tabular}
\caption{Key command-line arguments}
\end{table}

% =============================================================================
\section{Summary}
% =============================================================================

This system represents the first end-to-end, open-source pipeline
for automated riverine sand mining detection in India using:
\begin{itemize}[noitemsep]
  \item Sentinel-2 10\,m multispectral imagery
  \item 3--5 year quarterly historical trend features
  \item Six machine learning classifiers with cross-validated selection
  \item A human-in-the-loop labeling interface
  \item Interactive HTML probability maps
\end{itemize}

The Damodar Basin pilot (361\,km, 25 training samples, CV F1 = 0.82)
demonstrates technical feasibility. The architecture is explicitly
designed for national scaling to All-India coverage.

The key innovation --- combining \emph{current spectral features} with
\emph{multi-year temporal trend slopes} in a single ML feature vector
--- addresses the core limitation of single-date sand mining detection
and is grounded in recent literature from the Mekong, Ganga, and Ghana
river systems.

\vspace{1cm}
\begin{tcolorbox}[colback=lightblue!20, colframe=darkblue]
  \textbf{In plain language:} This system watches rivers from space for
  years, learns what sand mining looks like (both today and in terms of
  how the land has changed), and flags suspicious locations --- so
  inspectors know exactly where to go instead of guessing.
\end{tcolorbox}

\end{document}